% Pregunta 670 (2008-julio-OEP2008, n.º 22): condición de Solow. Función de esfuerzo en S
%   e(w) = 4,4/(1+exp(-1,8(w-2,8))) - 0,166  (e = 0 en w = 1)
%   Rayo desde el origen tangente en w* = 3,81, e* = 3,62 (pendiente e/w = 0,95, la máxima):
%   elasticidad del esfuerzo = 1 y coste por unidad de esfuerzo w/e mínimo.
% Autor: Víctor Gutiérrez Marcos
\documentclass[tikz,border=6pt]{standalone}
\input{tikz-tcee}
\begin{document}
\begin{tikzpicture}[tcee,x=1cm,y=1cm]
  \ejes{$w$}{$e$}{6}{5}
  \node[font=\small\bfseries,text=tceemorado,align=center] at (3,5.45)
    {Salario de eficiencia: condición de Solow};
  % Función de esfuerzo
  \draw[curva1,domain=1:5.6,samples=80] plot (\x,{4.4/(1+exp(-1.8*(\x-2.8)))-0.1658})
    node[etiqueta,right,text=tceemorado] {$e(w)$};
  % Rayo tangente desde el origen
  \draw[curva2] (0,0) -- (4.9,4.66) node[etiqueta,right,text=tceeazul] {rayo};
  \node[etiqueta,font=\footnotesize,align=left,text=tceeazul,anchor=south west] at (0.2,2.55) {rayo tangente:\\$e/w$ máxima};
  % Tangencia
  \proyeccion{(3.811,3.621)}{$w^*$}{$e^*$}
  \node[punto] at (3.811,3.621) {};
  \node[etiqueta,above left] at (3.811,3.621) {$E$};
  \node[etiqueta,font=\footnotesize,align=left,anchor=north west] at (3.9,2.6)
    {En $E$: $\dfrac{w\,e'(w)}{e(w)}=1$\\[2pt]$\Rightarrow\ w/e$ mínimo};
\end{tikzpicture}
\end{document}
